\section{Билет 59. Термодинамические процессы. Обратимые и необратимые процессы. Круговой процесс. Работа системы в циклическом процессе. Графический смысл. Тепловая машина. Определение и принципиальное устройство.}

\begin{definition}
\textbf{Термодинамическим процессом} называется переход системы из одного равновесного состояния в другое. В зависимости от возможности возврата в исходное состояние без изменений в окружающей среде процессы делятся на обратимые и необратимые.
\end{definition}

\begin{definition}
\textbf{Обратимым процессом} называется такой процесс, который может протекать как в прямом, так и в обратном направлении, причем при возвращении системы в исходное состояние в окружающей среде не происходит никаких изменений. Обратимые процессы являются идеализацией; они должны протекать бесконечно медленно (квазистатически) и без диссипации энергии (трения, теплопроводности и т.д.).
\par\noindent\textbf{Необратимые процессы} --- реальные процессы, после которых система не может самопроизвольно вернуться в исходное состояние без изменений во внешней среде. Все реальные процессы в природе необратимы из-за наличия диссипативных сил и конечной скорости протекания.
\end{definition}

\begin{definition}
\textbf{Круговым процессом} (или циклом) называется термодинамический процесс, в результате которого система возвращается в исходное состояние. На диаграмме $(p, V)$ равновесный цикл изображается замкнутой кривой.
\end{definition}

\begin{property}[Работа в цикле и её графический смысл]
Работа, совершаемая газом за полный цикл, численно равна площади фигуры, ограниченной кривой цикла на диаграмме $(p, V)$:
\begin{equation}
    A_{\text{цикл}} = \oint p \, dV.
\end{equation}
\begin{itemize}
    \item Если цикл обходится \textbf{по часовой стрелке} (прямой цикл), то работа газа за цикл положительна ($A > 0$). Такие циклы лежат в основе работы тепловых двигателей.
    \item Если цикл обходится \textbf{против часовой стрелки} (обратный цикл), работа газа отрицательна ($A < 0$). Такие циклы используются в холодильных машинах и тепловых насосах.
\end{itemize}
\end{property}

\begin{definition}
\textbf{Тепловой машиной} (тепловым двигателем) называется устройство, преобразующее внутреннюю энергию топлива в механическую работу.
\par\noindent\textbf{Принципиальное устройство:}
\begin{enumerate}
    \item \textbf{Нагреватель} --- тело с высокой температурой $T_1$, от которого рабочее тело получает количество теплоты $Q_1 > 0$.
    \item \textbf{Рабочее тело} --- газ или пар, совершающий круговой процесс.
    \item \textbf{Холодильник} --- тело с более низкой температурой $T_2 < T_1$, которому рабочее тело отдаёт часть энергии в виде теплоты $Q_2 < 0$ (или $|Q_2|$).
\end{enumerate}
За один цикл рабочее тело возвращается в исходное состояние ($\Delta U = 0$), поэтому согласно первому началу термодинамики полезная работа равна разности полученной и отданной теплоты:
\begin{equation}
    A = Q_1 - |Q_2|. \label{eq:work_cycle_59}
\end{equation}
\end{definition}
